\documentclass[12pt,a4paper]{article}

\usepackage[T2A]{fontenc}
\usepackage[utf8]{inputenc}
\usepackage[russian]{babel}
\usepackage{amsmath,amssymb,amsthm}
\usepackage{array}
\usepackage{enumitem}
\usepackage{listings}
\usepackage{xcolor}
\usepackage{geometry}

\geometry{left=25mm,right=20mm,top=20mm,bottom=20mm}

\lstset{
  basicstyle=\ttfamily\small,
  keywordstyle=\bfseries\color{blue!60!black},
  commentstyle=\itshape\color{green!40!black},
  stringstyle=\color{red!60!black},
  numbers=left,
  numberstyle=\tiny,
  stepnumber=1,
  numbersep=8pt,
  frame=single,
  breaklines=true,
  tabsize=2,
  showstringspaces=false,
  literate=
    {А}{{\CYRA}}1
    {Б}{{\CYRB}}1
    {В}{{\CYRV}}1
    {Г}{{\CYRG}}1
    {Д}{{\CYRD}}1
    {Е}{{\CYRE}}1
    {Ё}{{\CYRYO}}1
    {Ж}{{\CYRZH}}1
    {З}{{\CYRZ}}1
    {И}{{\CYRI}}1
    {Й}{{\CYRISHRT}}1
    {К}{{\CYRK}}1
    {Л}{{\CYRL}}1
    {М}{{\CYRM}}1
    {Н}{{\CYRN}}1
    {О}{{\CYRO}}1
    {П}{{\CYRP}}1
    {Р}{{\CYRR}}1
    {С}{{\CYRS}}1
    {Т}{{\CYRT}}1
    {У}{{\CYRU}}1
    {Ф}{{\CYRF}}1
    {Х}{{\CYRH}}1
    {Ц}{{\CYRC}}1
    {Ч}{{\CYRCH}}1
    {Ш}{{\CYRSH}}1
    {Щ}{{\CYRSHCH}}1
    {Ъ}{{\CYRHRDSN}}1
    {Ы}{{\CYRERY}}1
    {Ь}{{\CYRSFTSN}}1
    {Э}{{\CYREREV}}1
    {Ю}{{\CYRYU}}1
    {Я}{{\CYRYA}}1
    {а}{{\cyra}}1
    {б}{{\cyrb}}1
    {в}{{\cyrv}}1
    {г}{{\cyrg}}1
    {д}{{\cyrd}}1
    {е}{{\cyre}}1
    {ё}{{\cyryo}}1
    {ж}{{\cyrzh}}1
    {з}{{\cyrz}}1
    {и}{{\cyri}}1
    {й}{{\cyrishrt}}1
    {к}{{\cyrk}}1
    {л}{{\cyrl}}1
    {м}{{\cyrm}}1
    {н}{{\cyrn}}1
    {о}{{\cyro}}1
    {п}{{\cyrp}}1
    {р}{{\cyrr}}1
    {с}{{\cyrs}}1
    {т}{{\cyrt}}1
    {у}{{\cyru}}1
    {ф}{{\cyrf}}1
    {х}{{\cyrh}}1
    {ц}{{\cyrc}}1
    {ч}{{\cyrch}}1
    {ш}{{\cyrsh}}1
    {щ}{{\cyrshch}}1
    {ъ}{{\cyrhrdsn}}1
    {ы}{{\cyrery}}1
    {ь}{{\cyrsftsn}}1
    {э}{{\cyrerev}}1
    {ю}{{\cyryu}}1
    {я}{{\cyrya}}1
}

\newtheorem{definition}{Определение}
\newtheorem{statement}{Утверждение}
\newtheorem{example}{Пример}

\title{Билет №53 \\ \small Системы программирования (СП), типовые компоненты СП: языки, трансляторы, редакторы связей, отладчики, текстовые редакторы. Модульное программирование. Типы модулей. Связывание модулей по управлению и данным. (СпБГУ)}
\author{Сериков Владислав Александрович}
\date{16 мая 2026}

\begin{document}

\maketitle
\tableofcontents
\newpage

\section{Понятие системы программирования}

\begin{definition}
\textbf{Система программирования} --- это комплекс языковых, программных и инструментальных средств, предназначенных для разработки, трансляции, сборки, отладки, сопровождения и выполнения программ.
\end{definition}

Система программирования обычно включает:
\begin{enumerate}
  \item язык программирования или семейство языков;
  \item трансляторы: компиляторы, интерпретаторы, ассемблеры, препроцессоры;
  \item средства компоновки: редакторы связей, компоновщики, загрузчики;
  \item отладчики;
  \item текстовые редакторы или интегрированные среды разработки;
  \item библиотеки времени выполнения;
  \item системы управления проектом и сборкой;
  \item средства статического анализа, профилирования и тестирования.
\end{enumerate}

Главная задача системы программирования --- обеспечить переход от исходного текста программы к выполняемой программе, а также поддержать полный жизненный цикл разработки.

\section{Типовая схема обработки программы}

Разработка и выполнение программы обычно проходят следующие стадии:

\[
\text{исходный текст}
\longrightarrow
\text{препроцессор}
\longrightarrow
\text{компилятор}
\longrightarrow
\text{объектный модуль}
\longrightarrow
\text{редактор связей}
\longrightarrow
\text{исполняемый модуль}
\longrightarrow
\text{загрузчик}
\longrightarrow
\text{выполнение}
\]

Для языков с виртуальной машиной схема может выглядеть иначе:

\[
\text{исходный текст}
\longrightarrow
\text{компилятор в байт-код}
\longrightarrow
\text{байт-код}
\longrightarrow
\text{виртуальная машина/JIT-компилятор}
\longrightarrow
\text{выполнение}
\]

Для интерпретируемых языков возможна схема:

\[
\text{исходный текст}
\longrightarrow
\text{лексический и синтаксический анализ}
\longrightarrow
\text{интерпретация}
\]

На практике граница между компиляцией и интерпретацией часто размыта: современные языки могут использовать байт-код, JIT-компиляцию, кэширование промежуточного представления и динамическую оптимизацию.

\section{Языки программирования как компонент СП}

\subsection{Назначение языка программирования}

Язык программирования является формальным средством записи алгоритмов и структур данных. Он задаёт:
\begin{enumerate}
  \item лексические правила: допустимые символы, идентификаторы, ключевые слова;
  \item синтаксис: правила построения выражений, операторов, объявлений;
  \item семантику: смысл программных конструкций;
  \item систему типов;
  \item модель памяти;
  \item средства модульности;
  \item средства взаимодействия с внешними библиотеками и операционной системой.
\end{enumerate}

\begin{definition}
\textbf{Синтаксис языка} --- совокупность правил, определяющих допустимые последовательности символов и конструкций программы.
\end{definition}

\begin{definition}
\textbf{Семантика языка} --- совокупность правил, определяющих смысл синтаксически корректных программ.
\end{definition}

\subsection{Классификация языков программирования}

Языки программирования можно классифицировать по разным признакам.

\subsubsection{По уровню абстракции}

\begin{enumerate}
  \item \textbf{Машинные языки} --- программы записываются в машинных кодах.
  \item \textbf{Языки ассемблера} --- используют мнемоники машинных команд.
  \item \textbf{Языки высокого уровня} --- C, Pascal, Fortran, Java, Python, C++, Go и др.
  \item \textbf{Предметно-ориентированные языки} --- SQL, HTML, Verilog, MATLAB-подобные языки.
\end{enumerate}

\subsubsection{По парадигме программирования}

\begin{enumerate}
  \item \textbf{Императивные} --- программа задаёт последовательность команд.
  \item \textbf{Процедурные} --- основная единица декомпозиции: процедура или функция.
  \item \textbf{Объектно-ориентированные} --- программа строится из объектов и классов.
  \item \textbf{Функциональные} --- вычисления рассматриваются как применение функций.
  \item \textbf{Логические} --- программа описывает факты и правила вывода.
  \item \textbf{Декларативные} --- описывается, что нужно получить, а не как это вычислять.
\end{enumerate}

\subsubsection{По способу выполнения}

\begin{enumerate}
  \item \textbf{Компилируемые}: C, C++, Rust, Fortran.
  \item \textbf{Интерпретируемые}: классически Python, JavaScript, Ruby.
  \item \textbf{Языки с виртуальной машиной}: Java, C\#, Kotlin.
  \item \textbf{Смешанные модели}: JIT-компиляция, AOT-компиляция, байт-код.
\end{enumerate}

\section{Трансляторы}

\begin{definition}
\textbf{Транслятор} --- программа, переводящая программу с одного языка на другой с сохранением её смысла.
\end{definition}

Входной язык называют \textbf{исходным}, выходной --- \textbf{целевым}.

Основные виды трансляторов:
\begin{enumerate}
  \item компилятор;
  \item интерпретатор;
  \item ассемблер;
  \item препроцессор;
  \item декомпилятор;
  \item транспилятор.
\end{enumerate}

\subsection{Компилятор}

\begin{definition}
\textbf{Компилятор} --- транслятор, преобразующий программу с языка высокого уровня в эквивалентную программу на машинном языке, языке ассемблера, байт-коде или другом низкоуровневом представлении.
\end{definition}

Компилятор обычно не выполняет программу непосредственно, а порождает объектный или исполняемый код.

Типичные фазы компиляции:
\begin{enumerate}
  \item лексический анализ;
  \item синтаксический анализ;
  \item семантический анализ;
  \item построение промежуточного представления;
  \item оптимизация;
  \item генерация целевого кода;
  \item формирование объектного модуля.
\end{enumerate}

\subsection{Лексический анализ}

\begin{definition}
\textbf{Лексический анализ} --- фаза трансляции, на которой входной поток символов разбивается на лексемы.
\end{definition}

\begin{definition}
\textbf{Лексема} --- минимальная значимая последовательность символов языка: идентификатор, число, ключевое слово, знак операции, разделитель.
\end{definition}

Например, строка

\begin{lstlisting}[language=C]
x = a + 10;
\end{lstlisting}

разбивается на лексемы:
\[
\texttt{x},\quad \texttt{=},\quad \texttt{a},\quad \texttt{+},\quad \texttt{10},\quad \texttt{;}
\]

\subsection{Синтаксический анализ}

\begin{definition}
\textbf{Синтаксический анализ} --- проверка соответствия последовательности лексем грамматике языка и построение синтаксического дерева.
\end{definition}

Например, выражение

\[
a + b * c
\]

обычно разбирается как

\[
a + (b * c),
\]

поскольку умножение имеет более высокий приоритет, чем сложение.

\subsection{Семантический анализ}

\begin{definition}
\textbf{Семантический анализ} --- проверка смысловой корректности программы, не сводимой только к синтаксису.
\end{definition}

На этой стадии проверяются:
\begin{enumerate}
  \item объявлены ли все идентификаторы;
  \item согласованы ли типы;
  \item корректно ли используются функции;
  \item допустимы ли операции над объектами;
  \item соблюдены ли правила областей видимости;
  \item корректно ли выполняются преобразования типов.
\end{enumerate}

Пример синтаксически корректной, но семантически ошибочной программы:

\begin{lstlisting}[language=C]
int x;
x = "hello";
\end{lstlisting}

Синтаксис корректен, но типы несовместимы.

\subsection{Промежуточное представление}

\begin{definition}
\textbf{Промежуточное представление} --- внутренняя форма программы, удобная для анализа, оптимизации и генерации кода.
\end{definition}

Примеры промежуточных представлений:
\begin{enumerate}
  \item абстрактное синтаксическое дерево;
  \item трёхадресный код;
  \item байт-код;
  \item SSA-форма;
  \item граф потока управления.
\end{enumerate}

Пример трёхадресного кода для выражения

\[
x = a + b * c
\]

может иметь вид:

\begin{lstlisting}
t1 = b * c
t2 = a + t1
x = t2
\end{lstlisting}

\subsection{Оптимизация}

\begin{definition}
\textbf{Оптимизация программы} --- преобразование программы в эквивалентную по смыслу, но более эффективную по некоторому критерию форму.
\end{definition}

Критерии эффективности:
\begin{enumerate}
  \item время выполнения;
  \item объём машинного кода;
  \item потребление памяти;
  \item энергопотребление;
  \item количество обращений к памяти.
\end{enumerate}

Примеры оптимизаций:
\begin{enumerate}
  \item свёртка констант;
  \item удаление мёртвого кода;
  \item устранение общих подвыражений;
  \item инлайнинг функций;
  \item оптимизация циклов;
  \item распределение регистров.
\end{enumerate}

Пример свёртки констант:

\begin{lstlisting}[language=C]
x = 2 * 3 + y;
\end{lstlisting}

может быть преобразовано в

\begin{lstlisting}[language=C]
x = 6 + y;
\end{lstlisting}

\subsection{Генерация кода}

\begin{definition}
\textbf{Генерация кода} --- фаза компиляции, на которой промежуточное представление преобразуется в целевой код.
\end{definition}

Целевым кодом может быть:
\begin{enumerate}
  \item машинный код конкретного процессора;
  \item ассемблер;
  \item байт-код виртуальной машины;
  \item код на другом языке высокого уровня.
\end{enumerate}

\subsection{Интерпретатор}

\begin{definition}
\textbf{Интерпретатор} --- программа, которая непосредственно выполняет исходную программу или её промежуточное представление, не создавая заранее самостоятельный исполняемый файл.
\end{definition}

Преимущества интерпретации:
\begin{enumerate}
  \item удобство интерактивной работы;
  \item простота отладки;
  \item переносимость;
  \item возможность динамического изменения программы.
\end{enumerate}

Недостатки:
\begin{enumerate}
  \item обычно меньшая скорость выполнения;
  \item ошибки могут обнаруживаться позднее;
  \item зависимость от среды выполнения.
\end{enumerate}

\subsection{Ассемблер}

\begin{definition}
\textbf{Ассемблер} --- транслятор, переводящий программу с языка ассемблера в машинный код.
\end{definition}

Язык ассемблера является символическим представлением машинных команд.

Пример:

\begin{lstlisting}
mov eax, 1
add eax, 2
\end{lstlisting}

Ассемблер преобразует мнемоники команд и символические имена в машинные инструкции и адреса.

\subsection{Препроцессор}

\begin{definition}
\textbf{Препроцессор} --- программа, выполняющая предварительную обработку исходного текста до основной трансляции.
\end{definition}

Типичные функции препроцессора:
\begin{enumerate}
  \item подстановка макросов;
  \item включение файлов;
  \item условная компиляция;
  \item удаление комментариев;
  \item генерация кода по шаблонам.
\end{enumerate}

Пример на C:

\begin{lstlisting}[language=C]
#define N 10

int a[N];
\end{lstlisting}

После препроцессирования константа \texttt{N} может быть подставлена в текст программы.

\section{Объектные модули}

\begin{definition}
\textbf{Объектный модуль} --- результат раздельной трансляции части программы, содержащий машинный код, таблицу символов, информацию о внешних ссылках, данные для релокации и служебную информацию.
\end{definition}

Объектный модуль обычно ещё не является полностью исполняемой программой, поскольку может содержать обращения к внешним функциям и переменным.

Типичная структура объектного модуля:
\begin{enumerate}
  \item заголовок;
  \item секция кода;
  \item секция инициализированных данных;
  \item секция неинициализированных данных;
  \item таблица символов;
  \item таблица релокаций;
  \item отладочная информация.
\end{enumerate}

Например, файл \texttt{main.o}, полученный из \texttt{main.c}, может содержать вызов функции \texttt{sum}, но сама функция может находиться в другом объектном модуле \texttt{sum.o}.

\section{Редакторы связей и компоновщики}

\subsection{Основные определения}

\begin{definition}
\textbf{Редактор связей}, или \textbf{компоновщик}, --- программа, объединяющая объектные модули и библиотеки в исполняемый модуль или библиотеку.
\end{definition}

Основные задачи редактора связей:
\begin{enumerate}
  \item объединение объектных модулей;
  \item разрешение внешних ссылок;
  \item размещение секций в адресном пространстве;
  \item выполнение релокации;
  \item подключение библиотек;
  \item формирование исполняемого файла.
\end{enumerate}

\subsection{Внешние определения и внешние ссылки}

\begin{definition}
\textbf{Внешнее определение} --- символ, определённый в данном модуле и доступный другим модулям.
\end{definition}

\begin{definition}
\textbf{Внешняя ссылка} --- обращение из данного модуля к символу, определённому в другом модуле.
\end{definition}

Пример:

\begin{lstlisting}[language=C]
// main.c
int sum(int a, int b);

int main() {
    return sum(2, 3);
}
\end{lstlisting}

\begin{lstlisting}[language=C]
// sum.c
int sum(int a, int b) {
    return a + b;
}
\end{lstlisting}

Модуль \texttt{main.o} содержит внешнюю ссылку на \texttt{sum}. Модуль \texttt{sum.o} содержит внешнее определение \texttt{sum}. Редактор связей сопоставляет ссылку с определением.

\subsection{Статическая и динамическая компоновка}

\begin{definition}
\textbf{Статическая компоновка} --- связывание, при котором код библиотек включается в исполняемый файл на этапе сборки.
\end{definition}

Преимущества статической компоновки:
\begin{enumerate}
  \item исполняемый файл менее зависит от внешней среды;
  \item проще развёртывание;
  \item версии библиотек фиксируются на этапе сборки.
\end{enumerate}

Недостатки:
\begin{enumerate}
  \item увеличивается размер исполняемого файла;
  \item сложнее обновлять библиотечный код;
  \item несколько программ могут содержать дублирующиеся копии одной библиотеки.
\end{enumerate}

\begin{definition}
\textbf{Динамическая компоновка} --- связывание, при котором библиотечный код подключается во время загрузки или выполнения программы.
\end{definition}

Преимущества динамической компоновки:
\begin{enumerate}
  \item экономия памяти и дискового пространства;
  \item возможность обновления библиотек без полной пересборки программ;
  \item разделение библиотечного кода между процессами.
\end{enumerate}

Недостатки:
\begin{enumerate}
  \item зависимость от наличия совместимых библиотек;
  \item возможны конфликты версий;
  \item запуск программы может требовать дополнительного поиска и загрузки библиотек.
\end{enumerate}

\subsection{Релокация}

\begin{definition}
\textbf{Релокация} --- изменение адресов в объектном коде в соответствии с фактическим размещением модулей и секций в памяти или исполняемом файле.
\end{definition}

При раздельной компиляции компилятор не знает окончательные адреса всех функций и данных. Поэтому в объектном модуле оставляются специальные записи релокации.

Минимальный пример:

\begin{lstlisting}[language=C]
extern int x;

int f() {
    return x + 1;
}
\end{lstlisting}

Компилятор не знает адрес переменной \texttt{x}. Он создаёт машинный код с временным адресом и записью: ``здесь нужно подставить адрес символа \texttt{x}''. На этапе связывания компоновщик подставляет фактический адрес.

\section{Алгоритм работы редактора связей}

Типовой алгоритм статической компоновки можно описать следующим образом.

\subsection*{Вход}

\begin{enumerate}
  \item набор объектных модулей;
  \item набор библиотек;
  \item параметры компоновки;
  \item стартовый символ, например \texttt{main} или \texttt{\_start}.
\end{enumerate}

\subsection*{Выход}

Исполняемый модуль или библиотека.

\subsection*{Шаги алгоритма}

\begin{enumerate}
  \item Прочитать заголовки всех объектных модулей.
  \item Построить общую таблицу внешних символов.
  \item Для каждого символа определить, где он объявлен и где используется.
  \item Проверить наличие неразрешённых внешних ссылок.
  \item Проверить отсутствие недопустимых повторных определений.
  \item Разместить секции кода и данных в адресном пространстве.
  \item Вычислить окончательные адреса символов.
  \item Выполнить релокацию: исправить адресные поля машинных инструкций и данных.
  \item Подключить необходимые библиотечные модули.
  \item Сформировать заголовки исполняемого файла.
  \item Записать исполняемый модуль.
\end{enumerate}

\subsection*{Минимальный пример}

Пусть имеются два модуля:

\begin{lstlisting}[language=C]
// main.c
int inc(int);

int main() {
    return inc(5);
}
\end{lstlisting}

\begin{lstlisting}[language=C]
// inc.c
int inc(int x) {
    return x + 1;
}
\end{lstlisting}

После компиляции:

\[
\texttt{main.c} \rightarrow \texttt{main.o},
\qquad
\texttt{inc.c} \rightarrow \texttt{inc.o}.
\]

В \texttt{main.o} содержится:
\begin{enumerate}
  \item определение символа \texttt{main};
  \item внешняя ссылка на символ \texttt{inc}.
\end{enumerate}

В \texttt{inc.o} содержится:
\begin{enumerate}
  \item определение символа \texttt{inc}.
\end{enumerate}

Компоновщик строит таблицу:

\[
\begin{array}{|c|c|c|}
\hline
\text{Символ} & \text{Где определён} & \text{Где используется} \\
\hline
\texttt{main} & \texttt{main.o} & \text{стартовый код} \\
\hline
\texttt{inc} & \texttt{inc.o} & \texttt{main.o} \\
\hline
\end{array}
\]

Затем он заменяет обращение к \texttt{inc} в \texttt{main.o} на фактический адрес функции в итоговом исполняемом файле.

\section{Загрузчик}

\begin{definition}
\textbf{Загрузчик} --- компонент операционной системы, который размещает исполняемый модуль в памяти и подготавливает его к выполнению.
\end{definition}

Функции загрузчика:
\begin{enumerate}
  \item чтение исполняемого файла;
  \item создание адресного пространства процесса;
  \item загрузка сегментов кода и данных;
  \item настройка стека и кучи;
  \item подключение динамических библиотек;
  \item выполнение динамической релокации;
  \item передача управления начальной точке программы.
\end{enumerate}

\section{Отладчики}

\begin{definition}
\textbf{Отладчик} --- инструментальное средство, позволяющее исследовать выполнение программы, обнаруживать ошибки и анализировать состояние программы во время выполнения.
\end{definition}

Типичные возможности отладчика:
\begin{enumerate}
  \item запуск программы под управлением отладчика;
  \item пошаговое выполнение;
  \item установка точек останова;
  \item просмотр значений переменных;
  \item просмотр стека вызовов;
  \item изменение значений переменных во время выполнения;
  \item дизассемблирование машинного кода;
  \item анализ дампов памяти;
  \item отладка многопоточных программ.
\end{enumerate}

\subsection{Точки останова}

\begin{definition}
\textbf{Точка останова} --- условие, при выполнении которого отладчик приостанавливает выполнение программы.
\end{definition}

Виды точек останова:
\begin{enumerate}
  \item по адресу инструкции;
  \item по строке исходного текста;
  \item по входу в функцию;
  \item условная точка останова;
  \item точка наблюдения за изменением данных.
\end{enumerate}

\subsection{Пошаговое выполнение}

Различают:
\begin{enumerate}
  \item \textbf{step into} --- войти внутрь вызываемой функции;
  \item \textbf{step over} --- выполнить вызов функции как один шаг;
  \item \textbf{step out} --- выполнить до выхода из текущей функции;
  \item \textbf{continue} --- продолжить выполнение до следующей точки останова.
\end{enumerate}

\subsection{Минимальный пример отладки}

Программа:

\begin{lstlisting}[language=C]
#include <stdio.h>

int div2(int x) {
    return 10 / x;
}

int main() {
    int a = 0;
    printf("%d\n", div2(a));
    return 0;
}
\end{lstlisting}

Ошибка возникает при делении на ноль.

Типичный ход отладки:
\begin{enumerate}
  \item установить точку останова в функции \texttt{main};
  \item запустить программу;
  \item выполнить шаг до строки вызова \texttt{div2(a)};
  \item проверить значение \texttt{a};
  \item заметить, что \texttt{a = 0};
  \item войти в функцию \texttt{div2};
  \item убедиться, что выражение \texttt{10 / x} выполняется при \texttt{x = 0}.
\end{enumerate}

\section{Текстовые редакторы и интегрированные среды разработки}

\begin{definition}
\textbf{Текстовый редактор} --- программа для создания и изменения текстовых файлов, в том числе исходных текстов программ.
\end{definition}

Функции редактора, важные для программирования:
\begin{enumerate}
  \item подсветка синтаксиса;
  \item автодополнение;
  \item навигация по проекту;
  \item поиск и замена;
  \item работа с несколькими файлами;
  \item отображение ошибок компиляции;
  \item интеграция с системами контроля версий;
  \item форматирование кода.
\end{enumerate}

\begin{definition}
\textbf{Интегрированная среда разработки} --- программная среда, объединяющая редактор, средства сборки, отладчик, навигацию по коду, анализатор ошибок и другие инструменты разработки.
\end{definition}

IDE обычно включает:
\begin{enumerate}
  \item редактор кода;
  \item компилятор или интеграцию с ним;
  \item отладчик;
  \item систему управления проектами;
  \item средства рефакторинга;
  \item систему автодополнения;
  \item просмотр структуры программы;
  \item терминал или консоль вывода.
\end{enumerate}

\section{Модульное программирование}

\subsection{Основная идея}

\begin{definition}
\textbf{Модульное программирование} --- метод разработки программ, при котором программа разбивается на относительно независимые части, называемые модулями, каждая из которых имеет чётко определённый интерфейс и скрытую реализацию.
\end{definition}

Модульное программирование основано на следующих принципах:
\begin{enumerate}
  \item декомпозиция сложной задачи на подзадачи;
  \item сокрытие деталей реализации;
  \item минимизация связей между частями программы;
  \item повторное использование кода;
  \item независимая разработка и тестирование модулей;
  \item раздельная компиляция.
\end{enumerate}

\subsection{Модуль}

\begin{definition}
\textbf{Модуль} --- логически завершённая часть программы, имеющая собственное назначение, интерфейс и реализацию.
\end{definition}

Модуль обычно содержит:
\begin{enumerate}
  \item экспортируемые имена;
  \item импортируемые имена;
  \item локальные данные;
  \item локальные функции;
  \item код инициализации;
  \item возможно, код завершения.
\end{enumerate}

\subsection{Интерфейс и реализация}

\begin{definition}
\textbf{Интерфейс модуля} --- совокупность сущностей, доступных другим модулям: функций, типов, констант, классов, переменных.
\end{definition}

\begin{definition}
\textbf{Реализация модуля} --- внутренний код и данные, обеспечивающие выполнение функций модуля, но скрытые от других модулей.
\end{definition}

Пример на C:

\begin{lstlisting}[language=C]
// stack.h
#ifndef STACK_H
#define STACK_H

void push(int x);
int pop(void);
int empty(void);

#endif
\end{lstlisting}

\begin{lstlisting}[language=C]
// stack.c
#include "stack.h"

#define N 100

static int data[N];
static int top = 0;

void push(int x) {
    data[top++] = x;
}

int pop(void) {
    return data[--top];
}

int empty(void) {
    return top == 0;
}
\end{lstlisting}

Здесь файл \texttt{stack.h} задаёт интерфейс, а \texttt{stack.c} --- реализацию. Переменные \texttt{data} и \texttt{top} объявлены как \texttt{static}, поэтому они имеют внутреннее связывание и недоступны из других модулей.

\section{Типы модулей}

Модули можно классифицировать по нескольким основаниям.

\subsection{По роли в программе}

\begin{enumerate}
  \item \textbf{Главный модуль} --- содержит точку входа программы.
  \item \textbf{Функциональный модуль} --- реализует группу связанных функций.
  \item \textbf{Модуль данных} --- инкапсулирует структуры данных и операции над ними.
  \item \textbf{Сервисный модуль} --- предоставляет вспомогательные функции.
  \item \textbf{Интерфейсный модуль} --- обеспечивает взаимодействие с пользователем или внешней системой.
  \item \textbf{Библиотечный модуль} --- предназначен для повторного использования в разных программах.
\end{enumerate}

\subsection{По способу подключения}

\begin{enumerate}
  \item \textbf{Исходный модуль} --- файл с исходным кодом.
  \item \textbf{Объектный модуль} --- результат компиляции исходного модуля.
  \item \textbf{Библиотечный модуль} --- модуль, входящий в статическую или динамическую библиотеку.
  \item \textbf{Загружаемый модуль} --- подключается во время выполнения.
\end{enumerate}

\subsection{По видимости содержимого}

\begin{enumerate}
  \item \textbf{Открытый модуль} --- значительная часть содержимого доступна извне.
  \item \textbf{Закрытый модуль} --- доступ к внутренним данным ограничен интерфейсом.
  \item \textbf{Параметризованный модуль} --- поведение зависит от параметров типа, значения или конфигурации.
\end{enumerate}

\subsection{По состоянию}

\begin{enumerate}
  \item \textbf{Модуль без состояния} --- не хранит изменяемых внутренних данных между вызовами.
  \item \textbf{Модуль с состоянием} --- имеет внутренние переменные, влияющие на последующие вызовы.
\end{enumerate}

Пример модуля без состояния:

\begin{lstlisting}[language=C]
int max(int a, int b) {
    return a > b ? a : b;
}
\end{lstlisting}

Пример модуля с состоянием:

\begin{lstlisting}[language=C]
static int counter = 0;

int next_id(void) {
    return ++counter;
}
\end{lstlisting}

\section{Связность и сцепление модулей}

В модульном программировании важно различать две характеристики: связность и сцепление.

\begin{definition}
\textbf{Связность модуля} --- степень внутренней логической согласованности элементов одного модуля.
\end{definition}

Чем выше связность, тем лучше: модуль должен решать одну чёткую задачу.

\begin{definition}
\textbf{Сцепление модулей} --- степень зависимости одного модуля от другого.
\end{definition}

Чем ниже сцепление, тем лучше: изменение одного модуля не должно требовать изменения большого числа других модулей.

\subsection{Виды связности}

От плохой к хорошей:
\begin{enumerate}
  \item \textbf{Случайная связность} --- элементы объединены без ясной причины.
  \item \textbf{Логическая связность} --- элементы выполняют похожие, но разные функции.
  \item \textbf{Временная связность} --- элементы выполняются в один момент времени, например при инициализации.
  \item \textbf{Процедурная связность} --- элементы образуют последовательность действий.
  \item \textbf{Коммуникационная связность} --- элементы работают с одними и теми же данными.
  \item \textbf{Функциональная связность} --- все элементы модуля служат одной функции.
  \item \textbf{Информационная связность} --- модуль инкапсулирует структуру данных и операции над ней.
\end{enumerate}

\subsection{Виды сцепления}

От плохого к хорошему:
\begin{enumerate}
  \item \textbf{Сцепление по содержимому} --- один модуль обращается к внутренним частям другого.
  \item \textbf{Сцепление по общей области данных} --- модули используют общие глобальные данные.
  \item \textbf{Сцепление по управлению} --- один модуль передаёт другому управляющие флаги.
  \item \textbf{Сцепление по структуре данных} --- модуль передаёт другому сложную структуру, из которой используется только часть.
  \item \textbf{Сцепление по данным} --- модули взаимодействуют через параметры.
  \item \textbf{Сцепление по сообщениям} --- взаимодействие происходит через абстрактные сообщения или интерфейсы.
\end{enumerate}

\section{Связывание модулей по управлению}

\begin{definition}
\textbf{Связывание модулей по управлению} --- форма взаимодействия модулей, при которой один модуль передаёт управление другому или определяет порядок его выполнения.
\end{definition}

Основные формы связывания по управлению:
\begin{enumerate}
  \item вызов процедуры или функции;
  \item возврат из процедуры;
  \item передача функции обратного вызова;
  \item передача управляющего параметра;
  \item обработка событий;
  \item исключения;
  \item сопрограммы;
  \item сообщения в акторных системах.
\end{enumerate}

\subsection{Вызов процедуры}

Простейший случай связывания по управлению:

\begin{lstlisting}[language=C]
void print_hello(void) {
    printf("Hello\n");
}

int main(void) {
    print_hello();
    return 0;
}
\end{lstlisting}

Модуль, содержащий \texttt{main}, передаёт управление модулю, содержащему \texttt{print\_hello}. После завершения функции управление возвращается вызывающему модулю.

\subsection{Связывание через управляющий флаг}

Пример нежелательно сильного сцепления по управлению:

\begin{lstlisting}[language=C]
void process(int mode) {
    if (mode == 1) {
        // печать
    } else if (mode == 2) {
        // сохранение
    } else if (mode == 3) {
        // отправка
    }
}
\end{lstlisting}

Вызывающий модуль знает внутренние режимы работы функции. Это ухудшает модульность.

Более модульное решение --- разделить операции:

\begin{lstlisting}[language=C]
void print_data(void);
void save_data(void);
void send_data(void);
\end{lstlisting}

Или использовать таблицу функций:

\begin{lstlisting}[language=C]
typedef void (*Action)(void);

void run(Action action) {
    action();
}
\end{lstlisting}

\subsection{Обратные вызовы}

\begin{definition}
\textbf{Обратный вызов} --- функция, передаваемая одному модулю для последующего вызова из него.
\end{definition}

Пример:

\begin{lstlisting}[language=C]
#include <stdio.h>

typedef void (*Handler)(int);

void for_each(int *a, int n, Handler h) {
    for (int i = 0; i < n; ++i) {
        h(a[i]);
    }
}

void print_int(int x) {
    printf("%d\n", x);
}

int main(void) {
    int a[] = {1, 2, 3};
    for_each(a, 3, print_int);
    return 0;
}
\end{lstlisting}

Здесь модуль обхода массива управляет циклом, но конкретное действие задаётся внешней функцией \texttt{print\_int}.

\section{Связывание модулей по данным}

\begin{definition}
\textbf{Связывание модулей по данным} --- форма взаимодействия, при которой модули обмениваются значениями, структурами данных, ссылками, указателями, сообщениями или используют общие данные.
\end{definition}

Основные формы связывания по данным:
\begin{enumerate}
  \item передача параметров;
  \item возврат результата;
  \item использование структур данных;
  \item использование глобальных переменных;
  \item обмен сообщениями;
  \item работа через файлы, каналы, сокеты, базы данных;
  \item разделяемая память.
\end{enumerate}

\subsection{Передача параметров}

Наиболее предпочтительная форма взаимодействия:

\begin{lstlisting}[language=C]
int square(int x) {
    return x * x;
}
\end{lstlisting}

Функция получает данные через параметр и возвращает результат. Она не зависит от глобального состояния.

\subsection{Передача структуры}

\begin{lstlisting}[language=C]
typedef struct {
    int x;
    int y;
} Point;

int squared_distance(Point p) {
    return p.x * p.x + p.y * p.y;
}
\end{lstlisting}

Если функция использует всю структуру или значимую её часть, такое взаимодействие допустимо. Если же передаётся большая структура, а используется одно поле, сцепление возрастает.

\subsection{Глобальные данные}

Пример нежелательного связывания по общей области данных:

\begin{lstlisting}[language=C]
int current_user_id;

void module_a(void) {
    current_user_id = 10;
}

void module_b(void) {
    printf("%d\n", current_user_id);
}
\end{lstlisting}

Недостатки:
\begin{enumerate}
  \item трудно понять, кто изменяет данные;
  \item сложнее тестировать модули;
  \item возникают скрытые зависимости;
  \item повышается риск ошибок в многопоточных программах.
\end{enumerate}

Лучше передавать данные явно:

\begin{lstlisting}[language=C]
void print_user(int user_id) {
    printf("%d\n", user_id);
}
\end{lstlisting}

\section{Раздельная компиляция}

\begin{definition}
\textbf{Раздельная компиляция} --- организация сборки программы, при которой разные исходные модули компилируются независимо, а затем объединяются редактором связей.
\end{definition}

Преимущества раздельной компиляции:
\begin{enumerate}
  \item ускорение повторной сборки;
  \item независимая разработка модулей;
  \item сокрытие реализации;
  \item создание библиотек;
  \item уменьшение сложности больших программ.
\end{enumerate}

Пример:

\begin{lstlisting}
gcc -c main.c   # main.c -> main.o
gcc -c math.c   # math.c -> math.o
gcc main.o math.o -o app
\end{lstlisting}

Если изменился только файл \texttt{math.c}, достаточно перекомпилировать только его и затем повторно выполнить связывание:

\begin{lstlisting}
gcc -c math.c
gcc main.o math.o -o app
\end{lstlisting}

\section{Алгоритм раздельной сборки программы}

\subsection*{Вход}

\begin{enumerate}
  \item множество исходных модулей;
  \item зависимости между модулями;
  \item параметры компиляции;
  \item параметры компоновки.
\end{enumerate}

\subsection*{Выход}

Исполняемая программа или библиотека.

\subsection*{Шаги}

\begin{enumerate}
  \item Построить граф зависимостей исходных файлов.
  \item Определить, какие файлы были изменены.
  \item Для каждого изменённого исходного файла запустить компилятор.
  \item Получить объектные модули.
  \item Передать объектные модули и библиотеки редактору связей.
  \item Разрешить внешние ссылки.
  \item Сформировать исполняемый файл.
  \item При необходимости запустить тесты.
\end{enumerate}

\subsection*{Иллюстрация}

Пусть проект состоит из файлов:

\[
\texttt{main.c},\quad \texttt{stack.c},\quad \texttt{stack.h}.
\]

Зависимости:

\[
\texttt{main.c} \rightarrow \texttt{stack.h},
\qquad
\texttt{stack.c} \rightarrow \texttt{stack.h}.
\]

Если изменился \texttt{stack.h}, нужно перекомпилировать и \texttt{main.c}, и \texttt{stack.c}. Если изменился только \texttt{stack.c}, то \texttt{main.c} можно не компилировать заново.

\section{Инкапсуляция и сокрытие информации}

\begin{definition}
\textbf{Инкапсуляция} --- объединение данных и операций над ними в одном модуле с ограничением прямого доступа к внутреннему состоянию.
\end{definition}

\begin{definition}
\textbf{Сокрытие информации} --- принцип, согласно которому модуль должен скрывать детали реализации, предоставляя внешнему миру только необходимый интерфейс.
\end{definition}

Преимущества:
\begin{enumerate}
  \item реализацию можно менять без изменения клиентского кода;
  \item уменьшается вероятность некорректного использования данных;
  \item упрощается тестирование;
  \item повышается сопровождаемость.
\end{enumerate}

Пример: стек.

Плохой вариант:

\begin{lstlisting}[language=C]
int stack[100];
int top;
\end{lstlisting}

Любой модуль может изменить \texttt{top} напрямую.

Хороший вариант:

\begin{lstlisting}[language=C]
void push(int x);
int pop(void);
int empty(void);
\end{lstlisting}

Внутреннее представление скрыто.

\section{Интерфейс модуля как контракт}

Интерфейс модуля можно рассматривать как контракт между разработчиком модуля и его пользователями.

Контракт включает:
\begin{enumerate}
  \item сигнатуры функций;
  \item типы данных;
  \item предусловия;
  \item постусловия;
  \item инварианты;
  \item возможные ошибки;
  \item требования к ресурсам.
\end{enumerate}

\begin{definition}
\textbf{Предусловие} --- условие, которое должно быть истинным перед вызовом операции.
\end{definition}

\begin{definition}
\textbf{Постусловие} --- условие, которое должно быть истинным после корректного завершения операции.
\end{definition}

\begin{definition}
\textbf{Инвариант} --- условие, сохраняющееся во всех допустимых состояниях объекта или модуля.
\end{definition}

Пример для стека:
\begin{enumerate}
  \item предусловие \texttt{pop}: стек не пуст;
  \item постусловие \texttt{push}: размер стека увеличился на 1;
  \item инвариант: $0 \leq top \leq N$.
\end{enumerate}

\section{Минимальный пример модульной программы}

\subsection*{Файл \texttt{counter.h}}

\begin{lstlisting}[language=C]
#ifndef COUNTER_H
#define COUNTER_H

void counter_reset(void);
int counter_next(void);

#endif
\end{lstlisting}

\subsection*{Файл \texttt{counter.c}}

\begin{lstlisting}[language=C]
#include "counter.h"

static int value = 0;

void counter_reset(void) {
    value = 0;
}

int counter_next(void) {
    return ++value;
}
\end{lstlisting}

\subsection*{Файл \texttt{main.c}}

\begin{lstlisting}[language=C]
#include <stdio.h>
#include "counter.h"

int main(void) {
    counter_reset();

    printf("%d\n", counter_next());
    printf("%d\n", counter_next());

    return 0;
}
\end{lstlisting}

\subsection*{Сборка}

\begin{lstlisting}
gcc -c counter.c
gcc -c main.c
gcc counter.o main.o -o app
\end{lstlisting}

\subsection*{Анализ}

\begin{enumerate}
  \item \texttt{counter.h} задаёт интерфейс модуля.
  \item \texttt{counter.c} содержит реализацию.
  \item Переменная \texttt{value} скрыта с помощью \texttt{static}.
  \item \texttt{main.c} не знает, как именно хранится счётчик.
  \item Связывание по управлению происходит через вызовы \texttt{counter\_reset} и \texttt{counter\_next}.
  \item Связывание по данным происходит через возвращаемые значения функций.
\end{enumerate}

\section{Ошибки связывания}

При компоновке могут возникать ошибки.

\subsection{Неразрешённая внешняя ссылка}

Пример:

\begin{lstlisting}[language=C]
int f(void);

int main(void) {
    return f();
}
\end{lstlisting}

Если функция \texttt{f} нигде не определена, компоновщик выдаст ошибку вида:

\begin{lstlisting}
undefined reference to `f'
\end{lstlisting}

\subsection{Повторное определение символа}

Пример:

\begin{lstlisting}[language=C]
// a.c
int x = 1;
\end{lstlisting}

\begin{lstlisting}[language=C]
// b.c
int x = 2;
\end{lstlisting}

Если оба файла входят в одну программу, может возникнуть ошибка повторного определения символа \texttt{x}.

\subsection{Несовместимость интерфейса и реализации}

Пример:

\begin{lstlisting}[language=C]
// f.h
int f(int);
\end{lstlisting}

\begin{lstlisting}[language=C]
// f.c
double f(double x) {
    return x + 1;
}
\end{lstlisting}

Интерфейс и реализация не согласованы. В зависимости от языка и компилятора ошибка может быть обнаружена на этапе компиляции или привести к некорректной работе.

\section{Библиотеки}

\begin{definition}
\textbf{Библиотека} --- набор модулей, предназначенных для повторного использования.
\end{definition}

Виды библиотек:
\begin{enumerate}
  \item статические;
  \item динамические;
  \item системные;
  \item пользовательские;
  \item стандартные библиотеки языка;
  \item сторонние библиотеки.
\end{enumerate}

\subsection{Статическая библиотека}

Статическая библиотека содержит объектные модули, которые могут быть включены в исполняемый файл на этапе компоновки.

Условный пример:

\begin{lstlisting}
ar rcs libmath.a add.o mul.o
gcc main.o -L. -lmath -o app
\end{lstlisting}

\subsection{Динамическая библиотека}

Динамическая библиотека подключается при загрузке или во время выполнения.

Условный пример:

\begin{lstlisting}
gcc -fPIC -c math.c
gcc -shared math.o -o libmath.so
gcc main.o -L. -lmath -o app
\end{lstlisting}

\section{Модульность и области видимости}

\begin{definition}
\textbf{Область видимости} --- часть программы, в которой имя может быть использовано для обращения к соответствующей сущности.
\end{definition}

\begin{definition}
\textbf{Время жизни объекта} --- интервал выполнения программы, в течение которого объект существует в памяти.
\end{definition}

Для модульности особенно важны:
\begin{enumerate}
  \item локальные переменные;
  \item статические локальные переменные;
  \item глобальные переменные с внешним связыванием;
  \item глобальные переменные с внутренним связыванием;
  \item пространства имён;
  \item классы и пакеты.
\end{enumerate}

В C ключевое слово \texttt{static} на уровне файла ограничивает видимость имени текущим модулем трансляции:

\begin{lstlisting}[language=C]
static int helper(void) {
    return 42;
}
\end{lstlisting}

Функция \texttt{helper} не экспортируется для других модулей.

\section{Преимущества модульного программирования}

\begin{enumerate}
  \item \textbf{Управление сложностью}: большая задача разбивается на меньшие.
  \item \textbf{Повторное использование}: готовые модули можно применять в разных проектах.
  \item \textbf{Параллельная разработка}: разные разработчики работают над разными модулями.
  \item \textbf{Упрощение тестирования}: каждый модуль можно тестировать отдельно.
  \item \textbf{Упрощение сопровождения}: изменения локализуются.
  \item \textbf{Повышение надёжности}: скрытие реализации уменьшает число некорректных зависимостей.
  \item \textbf{Ускорение сборки}: при раздельной компиляции пересобираются только изменённые части.
\end{enumerate}

\section{Недостатки и трудности модульного подхода}

\begin{enumerate}
  \item необходимость проектирования интерфейсов;
  \item риск чрезмерной декомпозиции;
  \item накладные расходы на вызовы и абстракции;
  \item усложнение сборки;
  \item проблемы совместимости версий модулей;
  \item необходимость документирования контрактов;
  \item риск циклических зависимостей.
\end{enumerate}

\section{Циклические зависимости модулей}

\begin{definition}
\textbf{Циклическая зависимость} --- ситуация, при которой модуль $A$ зависит от модуля $B$, а модуль $B$ прямо или косвенно зависит от модуля $A$.
\end{definition}

Пример:

\[
A \rightarrow B,\qquad B \rightarrow A.
\]

Циклические зависимости нежелательны, потому что:
\begin{enumerate}
  \item усложняют понимание архитектуры;
  \item затрудняют раздельную компиляцию;
  \item мешают повторному использованию;
  \item усложняют тестирование;
  \item могут приводить к проблемам инициализации.
\end{enumerate}

Способы устранения:
\begin{enumerate}
  \item выделить общий интерфейс в третий модуль;
  \item применить инверсию зависимостей;
  \item использовать обратные вызовы;
  \item передавать зависимости явно через параметры;
  \item разделить слишком крупные модули.
\end{enumerate}

\section{Граф зависимостей модулей}

Модули программы удобно представлять ориентированным графом.

\begin{definition}
\textbf{Граф зависимостей модулей} --- ориентированный граф $G = (V, E)$, где вершины $V$ соответствуют модулям, а дуга $(A, B) \in E$ означает, что модуль $A$ зависит от модуля $B$.
\end{definition}

Если граф ацикличен, то модули можно упорядочить топологически и собирать в порядке зависимостей.

\subsection{Алгоритм топологической сортировки}

\subsection*{Назначение}

Топологическая сортировка позволяет упорядочить модули так, чтобы каждый модуль располагался после тех модулей, от которых он зависит.

\subsection*{Вход}

Ориентированный ациклический граф зависимостей $G = (V, E)$.

\subsection*{Выход}

Последовательность вершин, в которой для каждой дуги $(u, v)$ вершина $v$ идёт раньше $u$, если $u$ зависит от $v$.

\subsection*{Шаги алгоритма Кана}

\begin{enumerate}
  \item Для каждой вершины вычислить входную степень.
  \item Поместить в очередь все вершины с нулевой входной степенью.
  \item Пока очередь не пуста:
  \begin{enumerate}
    \item извлечь вершину $v$;
    \item добавить $v$ в результат;
    \item для каждой дуги $v \rightarrow w$ удалить эту дугу;
    \item уменьшить входную степень $w$;
    \item если входная степень $w$ стала равной нулю, добавить $w$ в очередь.
  \end{enumerate}
  \item Если в результате оказались не все вершины, то в графе есть цикл.
\end{enumerate}

\subsection*{Пример}

Пусть зависимости:

\[
\texttt{main} \rightarrow \texttt{stack},
\qquad
\texttt{stack} \rightarrow \texttt{memory},
\qquad
\texttt{main} \rightarrow \texttt{io}.
\]

Это означает, что:
\begin{enumerate}
  \item \texttt{main} зависит от \texttt{stack};
  \item \texttt{stack} зависит от \texttt{memory};
  \item \texttt{main} зависит от \texttt{io}.
\end{enumerate}

Возможный порядок сборки:

\[
\texttt{memory},\quad \texttt{io},\quad \texttt{stack},\quad \texttt{main}.
\]

\section{Практические рекомендации по проектированию модулей}

\begin{enumerate}
  \item Один модуль должен иметь одну основную ответственность.
  \item Интерфейс должен быть меньше реализации.
  \item Глобальные данные следует использовать только при строгой необходимости.
  \item Внутренние функции и данные нужно скрывать.
  \item Взаимодействие модулей лучше организовывать через параметры и возвращаемые значения.
  \item Не следует передавать управляющие флаги, если можно разделить операции.
  \item Следует избегать циклических зависимостей.
  \item Интерфейсы должны быть устойчивыми.
  \item Ошибки должны быть явно описаны в контракте модуля.
  \item Модуль должен быть удобен для автономного тестирования.
\end{enumerate}

\section{Итоговая схема системы программирования}

Систему программирования можно представить как совокупность взаимодействующих компонентов:

\[
\begin{array}{c}
\text{Язык программирования} \\
\downarrow \\
\text{Текстовый редактор / IDE} \\
\downarrow \\
\text{Препроцессор} \\
\downarrow \\
\text{Компилятор / интерпретатор / ассемблер} \\
\downarrow \\
\text{Объектные модули} \\
\downarrow \\
\text{Редактор связей} \\
\downarrow \\
\text{Исполняемый модуль} \\
\downarrow \\
\text{Загрузчик и среда выполнения} \\
\downarrow \\
\text{Отладчик, профилировщик, анализатор}
\end{array}
\]

\section{Краткие выводы}

\begin{enumerate}
  \item Система программирования --- это комплекс средств для создания, трансляции, сборки, отладки и выполнения программ.
  \item Основные компоненты СП: языки, трансляторы, редакторы связей, отладчики, текстовые редакторы и IDE.
  \item Компилятор переводит программу в объектный или исполняемый код, интерпретатор выполняет программу непосредственно или через промежуточное представление.
  \item Редактор связей объединяет объектные модули, разрешает внешние ссылки и выполняет релокацию.
  \item Отладчик позволяет исследовать выполнение программы и находить ошибки.
  \item Модульное программирование основано на разбиении программы на части с чёткими интерфейсами.
  \item Качественный модуль обладает высокой внутренней связностью и низким внешним сцеплением.
  \item Связывание модулей может происходить по управлению и по данным.
  \item Раздельная компиляция является важной технической основой модульного программирования.
  \item Хорошая модульная архитектура упрощает разработку, тестирование, сопровождение и повторное использование программ.
\end{enumerate}

\end{document}
